%% toc: P15 159 159.3,159.4,159.5,159.6
\chapter{Integration on the Conditioning Sigma-Algebra}
\label{sec:p15-w02c-p15-007-restricted-integrals}

The existence argument for conditional expectation in
\cref{thm:p15-conditional-existence} applies the Radon--Nikodym theorem
on the smaller measurable space \((\Omega,\mathcal G)\), where the
reference measure is the restriction of \(P\) to \(\mathcal G\).
The densities it produces are therefore characterized by integrals
against that restricted probability. The defining identity of
\cref{def:p15-general-conditional-expectation}, by contrast, is written
with integrals against \(P\) on \(\mathcal F_0\). This chapter proves the
identity that connects the two kinds of integral, and then carries the
density construction through to a version of \(E(X\mid\mathcal G)\) with
every integral taken against \(P\). It closes with a model in which the
restriction step fails, because restricting a sigma-finite measure to a
smaller sigma-algebra can destroy sigma-finiteness.

\emph{Inputs.} A probability space \((\Omega,\mathcal F_0,P)\), a
sub-sigma-algebra \(\mathcal G\subseteq\mathcal F_0\), and either a
\(\mathcal G\)-measurable extended real function or an integrable real
random variable \(X\).
\emph{Conclusions.} Integrals of \(\mathcal G\)-measurable functions over
events of \(\mathcal G\) have the same value, and the same existence and
integrability status, under \(P\) and under its restriction. The
densities of the two measures \(A\mapsto\int_AX^\pm\,dP\) on
\(\mathcal G\) are \(P\)-integrable, and their difference, made finite on
a \(\mathcal G\)-measurable null set, satisfies the \(P\)-integral identity
that defines a version of \(E(X\mid\mathcal G)\).

\section{The restricted probability}
\label{sec:p15-w02c-p15-007-restriction}

The notation follows \cref{def:p15-general-conditional-expectation}.
Durrett writes \(\mathcal F_o\) for the large sigma-algebra and
\(\mathcal F\) for the conditioning one \cite[p.~205]{DurrettPTEAuthor2019};
here they are \(\mathcal F_0\) and \(\mathcal G\).

\begin{hbdefinition}{def:p15-w02c-p15-007-restricted-probability}{Restriction to a sub-sigma-algebra}
Let \((\Omega,\mathcal F_0,P)\) be a probability space and
\(\mathcal G\subseteq\mathcal F_0\) a sigma-algebra. The
\emph{restriction} of \(P\) to \(\mathcal G\) is the set function
\[
                 P|_{\mathcal G}(A)\defeq P(A)\qquad(A\in\mathcal G).
\]
It is a probability measure on \((\Omega,\mathcal G)\): countable
additivity on \(\mathcal G\) is countable additivity of \(P\) on members
of \(\mathcal F_0\) that happen to lie in \(\mathcal G\), and
\(P|_{\mathcal G}(\Omega)=1\).

An extended real function \(W\) on \(\Omega\) is
\(\mathcal G\)-\emph{measurable} if \(\{W\in B\}\in\mathcal G\) for every
Borel set \(B\) of extended reals. Since \(\mathcal G\subseteq\mathcal F_0\),
such a \(W\) is also \(\mathcal F_0\)-measurable, so both
\(\int W\,dP\) and \(\int W\,d(P|_{\mathcal G})\) are governed by
\cref{def:p15-nonnegative-integral} and
\cref{def:p15-integrable-functions}, on the measure spaces
\((\Omega,\mathcal F_0,P)\) and \((\Omega,\mathcal G,P|_{\mathcal G})\)
respectively. For \(A\in\mathcal G\) the function \(W\mathbf1_A\) is again
\(\mathcal G\)-measurable, so \(\int_AW\,d(P|_{\mathcal G})\) is meaningful
under the same conventions.

An \(\mathcal F_0\)-measurable function that is not \(\mathcal G\)-measurable
has no integral against \(P|_{\mathcal G}\): the lower simple functions in
\cref{def:p15-nonnegative-integral} must be measurable for the space on
which the integral is taken.
\index{restriction of a measure}\index{sub-sigma-algebra!restricted probability}
\end{hbdefinition}

\begin{hblemma}{thm:p15-w02c-p15-007-restricted-integral}{proved=annex:p15-w02c-p15-007-restricted}{Integration against the restricted probability}
In the setting of \cref{def:p15-w02c-p15-007-restricted-probability}, let
\(W\) be a \(\mathcal G\)-measurable extended real function and let
\(A\in\mathcal G\).
\begin{enumerate}
\item If \(W\geq0\), then
      \[
        \int_AW\,d(P|_{\mathcal G})=\int_AW\,dP,
      \]
      as an equality in \([0,\infty]\).
\item \(\int|W|\,d(P|_{\mathcal G})=E|W|\). Hence \(W\) is
      \(P|_{\mathcal G}\)-integrable exactly when it is \(P\)-integrable.
      The signed integral of \(W\) over \(A\) exists under one of the two
      measures exactly when it exists under the other, and the two values
      are then equal:
      \[
        \int_AW\,d(P|_{\mathcal G})=\int_AW\,dP.
      \]
\item For \(\mathcal G\)-measurable extended real \(V,W\), the event
      \(\{V>W\}\) belongs to \(\mathcal G\). Consequently
      \(V\leq W\) holds \(P|_{\mathcal G}\)-almost everywhere exactly when
      it holds \(P\)-almost everywhere, and the same is true of
      \(V=W\).
\end{enumerate}
\end{hblemma}

Item~3 matters because the two notions of null set differ in general:
a \(P\)-null event of \(\mathcal F_0\) outside \(\mathcal G\) is not in the
domain of \(P|_{\mathcal G}\). For comparisons between
\(\mathcal G\)-measurable functions the exceptional set is itself in
\(\mathcal G\), so the difference disappears. Thus the
\(P|_{\mathcal G}\)-almost-everywhere uniqueness of Radon--Nikodym
densities and the \(P\)-almost-everywhere identification of versions in
\cref{def:p15-general-conditional-expectation} describe the same classes
of \(\mathcal G\)-measurable functions.

\section{The density construction against \(P\)}
\label{sec:p15-w02c-p15-007-construction}

\begin{hbtheorem}{thm:p15-w02c-p15-007-density-version}{proved=annex:p15-w02c-p15-007-density}{The Radon--Nikodym densities give a version of the conditional expectation}
Let \(X\) be an integrable real random variable on
\((\Omega,\mathcal F_0,P)\) and let \(\mathcal G\subseteq\mathcal F_0\) be
a sigma-algebra. For \(A\in\mathcal G\) put
\[
            \nu_+(A)\defeq\int_AX^+\,dP,\qquad
            \nu_-(A)\defeq\int_AX^-\,dP.
\]
\begin{enumerate}
\item \(\nu_+\) and \(\nu_-\) are finite measures on
      \((\Omega,\mathcal G)\), each absolutely continuous with respect to
      \(P|_{\mathcal G}\). Hence, by \cref{thm:p15-radon-nikodym} on
      \((\Omega,\mathcal G)\), there are \(\mathcal G\)-measurable
      \(Z_\pm:\Omega\to[0,\infty]\) with
      \(\nu_\pm(A)=\int_AZ_\pm\,d(P|_{\mathcal G})\) for all
      \(A\in\mathcal G\).
\item Any such \(Z_\pm\) satisfy, for every \(A\in\mathcal G\),
      \[
          \int_AZ_\pm\,dP=\nu_\pm(A),\qquad E Z_\pm=EX^\pm<\infty.
      \]
\item The event \(N\defeq\{Z_+=\infty\}\cup\{Z_-=\infty\}\) belongs to
      \(\mathcal G\) and \(P(N)=0\).
\item The function \(Z\defeq(Z_+-Z_-)\mathbf1_{N^c}\), equal to zero on
      \(N\), is real valued, \(\mathcal G\)-measurable and
      \(P\)-integrable, and
      \[
          \int_AZ\,dP=\int_AX\,dP\qquad(A\in\mathcal G).
      \]
      Thus \(Z\) is a version of \(E(X\mid\mathcal G)\) in the sense of
      \cref{def:p15-general-conditional-expectation}.
\end{enumerate}
No completion of \(\mathcal G\) is used, and the null set \(N\) on which
the representatives are changed lies in \(\mathcal G\).
\end{hbtheorem}

This is the construction of \cref{thm:p15-conditional-existence}, with
the passage from \(P|_{\mathcal G}\)-integrals to \(P\)-integrals proved by
\cref{thm:p15-w02c-p15-007-restricted-integral} rather than assumed.
Durrett's argument takes the reference measure to be \(P\) on the smaller
sigma-algebra and writes the density's event integrals against \(P\)
directly \cite[p.~207]{DurrettPTEAuthor2019}; item~2 is the step that
justifies doing so. Durrett treats \(X\geq0\) first and then subtracts the
conditional expectations of \(X^+\) and \(X^-\); item~4 performs the same
subtraction on the two densities.

\begin{hbexample}{def:p15-w02c-p15-007-two-cell}{A two-cell conditioning sigma-algebra}
Let \(\Omega=[0,1]\), let \(\mathcal F_0\) be its Borel sigma-algebra and
let \(P\) be Lebesgue measure. Put \(L=[0,1/2)\), \(R=[1/2,1]\) and
\(\mathcal G=\{\emptyset,L,R,\Omega\}\). A function is
\(\mathcal G\)-measurable exactly when it is constant on \(L\) and constant
on \(R\): a function taking two values on \(L\) would have a level set
meeting \(L\) in a nonempty proper subset, and no member of \(\mathcal G\)
is of that form. The restricted probability gives \(L\) and \(R\) mass
\(1/2\) each, and assigns no mass to any other set, because no other set
is in its domain.

Take \(X(\omega)=\omega\). Then \(X=X^+\), \(X^-=0\), and
\[
        \nu_+(L)=\int_0^{1/2}\omega\,d\omega=\frac18,\qquad
        \nu_+(R)=\int_{1/2}^{1}\omega\,d\omega=\frac38.
\]
A \(\mathcal G\)-measurable density is \(Z_+=a\mathbf1_L+b\mathbf1_R\).
Its integral over \(L\) against \(P|_{\mathcal G}\) is the simple integral
\(a\,P|_{\mathcal G}(L)=a/2\), and over \(R\) it is \(b/2\). Matching
\(\nu_+\) gives \(a=1/4\) and \(b=3/4\); the values on the empty set and
on \(\Omega\) then match by additivity. Here \(Z_-=0\) and \(N=\emptyset\),
so \(Z=\tfrac14\mathbf1_L+\tfrac34\mathbf1_R\).

The \(P\)-integral identity can be checked directly:
\(\int_LZ\,dP=\tfrac14\cdot\tfrac12=\tfrac18=\int_LX\,dP\), and likewise
on \(R\). In this example both integrals of \(Z\) reduce to the same finite
sum because \(Z\) is a \(\mathcal G\)-simple function, which is the first
step of the proof of \cref{thm:p15-w02c-p15-007-restricted-integral}.

The hypothesis of \(\mathcal G\)-measurability cannot be dropped from that
lemma. The variable \(X\) itself is \(P\)-integrable with
\(\int_LX\,dP=1/8\), but \(\{X<1/4\}=[0,1/4)\notin\mathcal G\), so
\(X\) is not \(\mathcal G\)-measurable and \(\int X\,d(P|_{\mathcal G})\)
is not defined. Conditional expectation supplies the
\(\mathcal G\)-measurable replacement \(Z\) whose \(P\)-integrals over
\(\mathcal G\) agree with those of \(X\).
\index{conditional expectation!two-cell example}
\end{hbexample}

\section{Where restriction fails}
\label{sec:p15-w02c-p15-007-failure}

For a probability space the restriction is again a probability, so the
sigma-finiteness hypothesis of \cref{thm:p15-radon-nikodym} holds on
\((\Omega,\mathcal G)\) automatically. For an infinite measure it can
fail, and then the density step of the construction has no output.

\begin{hbtheorem}{thm:p15-w02c-p15-007-restriction-sigma-finite}{proved=annex:p15-w02c-p15-007-sigma-finite}{Restriction and sigma-finiteness}
Let \(\lambda\) be Lebesgue measure on the Borel sets of \(\mathbb R\),
let \(\mathcal G=\{\emptyset,\mathbb R\}\), and let
\(\nu(A)=\int_A\mathbf1_{[0,1]}\,d\lambda\) for Borel \(A\).
Write \(\lambda|_{\mathcal G}\) and \(\nu|_{\mathcal G}\) for the set
functions that agree with \(\lambda\) and \(\nu\) on \(\mathcal G\); as in
\cref{def:p15-w02c-p15-007-restricted-probability}, they are measures on
\((\mathbb R,\mathcal G)\) because countable additivity is inherited.
\begin{enumerate}
\item \(\lambda\) is sigma-finite, but \(\lambda|_{\mathcal G}\) is not.
\item \(\nu|_{\mathcal G}\) is a finite measure with
      \(\nu|_{\mathcal G}\ll\lambda|_{\mathcal G}\).
\item No \(\mathcal G\)-measurable \(g:\mathbb R\to[0,\infty]\) satisfies
      \(\nu(\mathbb R)=\int g\,d(\lambda|_{\mathcal G})\).
\end{enumerate}
\end{hbtheorem}

In this model the conclusion of \cref{thm:p15-radon-nikodym} fails
because the reference measure is not sigma-finite on the space where the
density is sought, although it is sigma-finite on a larger space. The construction of
\cref{thm:p15-w02c-p15-007-density-version} therefore uses the finiteness
of \(P\) in an essential way at the restriction step.
\Cref{thm:p15-w02c-p15-007-restricted-integral} is stated and proved here
for a probability measure, which is the case that construction needs.

\section{Relationships and sources}
\label{sec:p15-w02c-p15-007-relationships}

\begin{itemize}
\item \Cref{thm:p15-conditional-existence} states existence,
      integrability, the bound \(E|C_{\mathcal G}X|\leq E|X|\) and
      almost-everywhere uniqueness; its annex proof applies the density
      construction, and \cref{thm:p15-w02c-p15-007-density-version}
      writes every integral of that construction against \(P\).
\item \Cref{thm:p15-expectation-law-integral} gives a second derivation of
      \cref{thm:p15-w02c-p15-007-restricted-integral}: the identity map
      \((\Omega,\mathcal F_0)\to(\Omega,\mathcal G)\) is a random element
      whose law is \(P|_{\mathcal G}\). The annex records this derivation
      after the direct proof.
\item \Cref{thm:p15-measurable-limits} and \cref{thm:p15-monotone-integral}
      supply the approximation and the limit used in the direct proof.
\item \Cref{thm:p15-radon-nikodym} is applied on \((\Omega,\mathcal G)\),
      and \cref{thm:p15-w02c-p15-007-restriction-sigma-finite} shows why its
      sigma-finiteness hypothesis must be checked on that space.
\end{itemize}

Source: Durrett, \emph{Probability: Theory and Examples}, author version
of 11 January 2019, Section~4.1: the definition of conditional
expectation \cite[pp.~205--206]{DurrettPTEAuthor2019} and its existence
proof by the Radon--Nikodym theorem \cite[p.~207]{DurrettPTEAuthor2019}.
